Joint continuous-parameter inference provides a second route to separating
signal and noise. We evaluate a Python implementation of VALSE Heuristic~2
from the paper equations \cite{badiu2017valse}, retaining frequency and gain
uncertainty in the noise update. VALSE learns from one target block over a
circular frequency support, while the refined control uses the delay bound
$[0,12]$ and 16 calibration blocks for its residual statistic. The comparison
therefore measures pilot-domain reconstruction and inferred noise under these
different resources. Full-channel interpolation would additionally require a
physical delay branch for the frequency posterior.

Three new seeds give 450 paired groups for each of three sparse profiles
at 10 and 20~dB. The primary setting uses a 500-iteration limit and a relative
reconstruction-change tolerance of $10^{-6}$. All 2700 outputs are retained:
2688 met the stopping criterion, and 12 reached the iteration limit.
Appendix~\ref{app:valse} details the implementation and numerical checks.
Replaying the 12 capped inputs with a 5000-iteration ceiling gives 11 further
convergences. Across the six cells, pilot NMSE changes by at most
0.000624~dB and the mean inferred-noise ratio by $2.72\times10^{-5}$.
The primary table retains the original 500-iteration outputs.

\begin{table}[t]
\centering
\caption{Pilot-domain reconstruction and noise inference; 450 groups per cell. Columns 3--6 give NMSE in dB. VALSE learns noise from one block with circular frequency support; refined controls use the delay bound $[0,12]$. The comparison concerns the observed pilot grid.}
\label{tab:valse}
\small\setlength{\tabcolsep}{4pt}
\begin{tabular}{lrrrrrr}
\toprule
& SNR & Integer & Refined & Refined & VALSE & VALSE \\
Profile & (dB) & EB & residual & known $v$ & & $\mathbb E[\widehat v]/v$ \\
\midrule
Sparse integer & 10 & -14.86 & -14.44 & -14.55 & -14.30 & 1.04 \\
Sparse integer & 20 & -24.98 & -24.86 & -24.86 & -23.86 & 1.13 \\
Separated fractional & 10 & -10.95 & -14.09 & -14.57 & -14.46 & 1.04 \\
Separated fractional & 20 & -13.76 & -22.13 & -24.97 & -23.92 & 1.17 \\
Close fractional & 10 & -10.75 & -12.76 & -13.24 & -12.97 & 1.12 \\
Close fractional & 20 & -13.82 & -18.92 & -23.83 & -19.85 & 1.94 \\
\bottomrule
\end{tabular}
\end{table}

Table~\ref{tab:valse} shows noise inference with a continuous probabilistic
model. At 20~dB, the mean noise ratios were 1.13, 1.17, and 1.94 for integer,
separated fractional, and close fractional paths, respectively. The larger
close-pair ratio accompanies a larger pilot reconstruction error, identifying
a condition where the specified joint inference retains appreciable calibration
error. For separated fractional paths, pilot NMSE was $-23.92$~dB,
between the residual-calibrated and known-noise refined controls. These
pilot-domain observations complement the full-channel CP intervention by
showing how calibration also depends on the inferential model.
